%MIT OpenCourseWare: https://ocw.mit.edu
%RES.18-012 Algebra II Student Notes, Spring 2022
%License: Creative Commons BY-NC-SA 
%For information about citing these materials or our Terms of Use, visit: https://ocw.mit.edu/terms.

% \rhead{\rightmark}
\lhead{Lecture 19: Modules over a Ring}

\section{Modules over a Ring}

The motivation for modules is that we are trying to tell a story where rings are the protagonist, and for a story to be interesting, the protagonist must act. When we find a way for a ring to act, we get the definition of a module. 

\begin{definition}
    Let $R$ be a ring. A \textbf{module} $M$ over $R$ is an abelian group, together with an \textbf{action map} $R \times M \to M$ (written as $(r, m) \mapsto r(m)$ or $rm$), subject to the following axioms:
    \begin{itemize}
        \item $1_R(m) = m$ for all $m \im M$;
        \item $r_1(r_2(m)) = (r_1r_2)(m)$ for all $r_1, r_2 \in R$ and $m \in M$;
        \item Distributivity in both variables: $(r_1 + r_2)m = r_1m + r_2m$ for all $r_1, r_2 \in R$ and $m \in M$, and $r(m_1 + m_2) = rm_1 + rm_2$ for all $r \in R$ and $m_1, m_2 \in M$. 
    \end{itemize}
\end{definition}

The first two axioms are very similar to the definition of a group action on a set. So a ring to a module is like a group $G$ to a $G$-set (a set with an action by $G$). It's not exactly the same because in a ring we have two operations instead of one, but they're a similar flavor.

\subsection{Examples}

\vspace{0.75em}

\begin{example}
    If $R = F$ is a field, then a module is the same as a vector space. 
\end{example}

The axioms here are exactly the same. The textbook emphasizes this heavily, and this analogy can get you some mileage; but for general rings, things become more complicated. 

\begin{note}
    The definition also applies to a \emph{noncommutative} ring $R$, in the same way --- our definition does not reference commutativity. Then we have some familiar examples of modules over noncommutative rings: for example, given any field $F$ we can take $R = \operatorname{Mat}_{n \times n}(F)$ and $M = F^n$, since matrices act on column vectors by multiplication. As another example, if $R = \mathbb{C}[G]$ is the group ring, then a $R$-module is the same as a complex representation of $G$. 
\end{note}

For any ring $R$, there is a uniquely defined homomorphism $\mathbb{Z} \to R$, where $1 \mapsto 1_R$. On a similar note, every abelian group has a unique structure of a $\mathbb{Z}$-module: we know that $1$ (in $\mathbb{Z}$) must map $m \mapsto m$, so then by distributivity, $n = 1 + 1 + \cdots + 1$ must map \[v \mapsto \underbrace{v + v + \cdots + v}_n.\] Similarly, $-n$ must map $v$ to $-(v + v + \cdots + v)$. So a $\mathbb{Z}$-module is the same as an abelian group. 

\begin{example}
    What is a module over $R = \mathbb{C}[x]$?
\end{example}

\begin{proof}
    First, a $\CC[x]$-module is a $\mathbb{C}$-vector space $V$ by looking at the action of constant polynomials (which are just scalars). But then we also need to see what $x$ does. We know $x$ must act by a linear map $A : V \to V$, where $xv = Av$. There are no constraints on this map, and this defines the action of every other polynomial: so a $R$-module is a vector space $V$, together with a linear map $A : V \to V$. Explicitly, the action of a general polynomial $P(x) = a_nx^n + \cdots + a_0$ is given by \[Pv = a_0v + a_1Av + \cdots + a_nA^nv.\]  Note that the vector space may or may not be finite-dimensional; if it is, then we end up in a situation studied in linear algebra, where we have a vector space and a linear operator.
\end{proof}

\begin{example}
    What is a module over $R = \ZZ/n\ZZ$?
\end{example}

\begin{proof}
    The main point is that if $R/I$ is a quotient of $R$, then every $R/I$-module is also a $R$-module, where we define $r(m)$ to be $\overline{r}(m)$ (here $\overline{r}$ denotes $r$ mod $I$). Meanwhile, in order to go backwards, $I$ must act by $0$. So a $R/I$ module is the same as a $R$-module where every element of $I$ acts in a trivial way (meaning that $rv = 0$ for all $r \in I$ and $v \in M$). 

    So in this case, a $\ZZ/n\ZZ$-module is the same as an abelian group where the order of every element divides $n$ --- meaning $na = 0$ for all $a$ in the group. 

    Then more concretely, for every $m$ (where we use $\overline{m}$ to denote $m$ mod $n$), we can write \[\overline{m}v = \underbrace{v + v + \cdots + v}_m.\] In order for this to be well-defined, the sum should not depend on the choice of representative for the residue; but this is guaranteed by the condition $na = 0$. (This is the same reasoning as in the first paragraph, for this specific example.)
\end{proof}

For any ring $R$, there is a simple example of a module:

\begin{definition}
The \textbf{free module} over $R$ is $M = R$ itself, where the action is multiplication (meaning that $r(x) = rx$).
\end{definition} 

This is parallel to the observation that a group $G$ acts on itself by left multiplication. 

\subsection{Submodules}

\vspace{0.75em}

\begin{definition}
    Given a module $M$, a \textbf{submodule} $N \subset M$ is an abelian subgroup which is invariant under the $R$-action --- meaning $rx \in N$ for all $x \in N$ and $r \in R$. 
\end{definition}

If $N \subset M$ is a submodule, we can define their \textbf{quotient} $M/N$, where we take the quotient in the sense of abelian groups. This quotient of abelian groups carries a module structure as well, given by the obvious rule $r\overline{m} = \overline{rm}$ (where $\overline{m}$ denotes $m$ mod $N$). This is well-defined because $N$ is a submodule --- if $m_1 - m_2$ is in $N$, then $rm_1 - rm_2 = r(m_1 - r_2)$ is in $N$ as well.

Then the homomorphism theorem and correspondence theorem work in the exact same way as in abelian groups. (For rings and ideals, we saw they work in a similar way; but here the parallel is closer.)

\begin{example}
    What are the submodules of the free module $M = R$?
\end{example}

\begin{proof}
    The answer is exactly the ideals of $R$ --- we're looking for abelian subgroups of $R$ which are invariant under multiplication by all terms in $R$, and by definition these are ideals. 
\end{proof}

We'll later see how to understand \emph{any} module by looking at generators and relations --- this turns out to be easier than the corresponding problem for a group. But first we'll look at another example of a module, which will be useful for developing that theory. 

\begin{definition}
    Given two modules $M$ and $N$, their \textbf{direct sum} is \[M \oplus N = \{(m, n) \mid m \in M, n \in N\}\] with the action \[r(m, n) = (rm, rn).\] 
\end{definition}

\begin{note}
    The direct sum is the same as the product $M \times N$. This is true for any \emph{finite} sum --- we have \[M_1 \oplus \cdots \oplus M_n = M_1 \times \cdots \times M_n.\] But this isn't true for \emph{infinite} sums and products. 
\end{note}

\begin{definition}
    The \textbf{free module} of \textbf{rank} $n$ is \[R^n = \underbrace{R \oplus R \oplus \cdots \oplus R}_n.\] 
\end{definition}

In the case where $R = F$ is a field, the free module of rank $n$ is exactly $F^n$, the standard $n$-dimensional vector space. 

\subsection{Homomorphisms}

In linear algebra, we work with matrices in order to understand linear maps. Matrices are also relevant here --- the terms are different, but the concept is very similar. 

\begin{definition}
    A \textbf{homomorphism} from a module $M$ to a module $N$ is a homomorphism of abelian groups $\varphi : M \to N$, which is compatible with the $R$-action --- meaning $\varphi(rm) = r\varphi(m)$ for all $r \in R$ and $m \in M$.
\end{definition}

In vector spaces, this is the same as a linear map. 

We'll use $\operatorname{Hom}_R(M, N)$ to denote the set of all homomorphisms $M \to N$. Note that homomorphisms can be added and rescaled, in the same way as linear maps: $(\varphi_1 + \varphi_2)(m) = \varphi_1(m) + \varphi_2(m)$, and $(r\varphi)(m) = r\varphi(m)$. So then $\operatorname{Hom}_R(M, N)$ is itself a $R$-module. 

Understanding homomorphisms in general may be hard, but it's easy to understand homomorphisms from a free module. Given a homomorphism $\varphi \in \operatorname{Hom}_R(R, M)$ for any module $M$, we can let $m = \varphi(1_R)$. Then this determines the entire homomorphism --- for any $r \in R$, we have \[\varphi(r) = \varphi(r\cdot 1_R) = r\cdot \varphi(1_R) = rm.\] 

So a homomorphism is determined by $m = \varphi(1_R)$, and there are no restrictions on $m$ --- this is why $R$ is called a free module. This means $\operatorname{Hom}_R(R, M)$ is isomorphic to $M$: more explicitly, the bijection is given by mapping $\varphi \in \operatorname{Hom}_R(R, M)$ to $m_\varphi = \varphi(1)$, and $m \in M$ to the homomorphism $\varphi_m : r \mapsto rm$. 

Similarly, $\operatorname{Hom}_R(R^n, M)$ is equally easy to understand. Now $R^n$ is generated by the elements $1_i$ which have a $1$ in their $i$th place, and $0$'s everywhere else (so $1_i = (0, \ldots, 0, 1, 0, \ldots, 0)$, where the $1$ is in the $i$th place). So $\operatorname{Hom}_R(R^n, M)$ is isomorphic to $M^n$, where the bijection sends $\varphi \in \operatorname{Hom}_R(R^n, M)$ to the element $(\varphi(1_1), \varphi(1_2), \ldots, \varphi(1_n))$, and $(m_1, \ldots, m_n) \in M$ to the homomorphism $\varphi(x_1, \ldots, x_n) = \sum x_im_i$. 

In particular, we have $\operatorname{Hom}_R(R^n, R^m) = (R^m)^n = \operatorname{Mat}_{m \times n}(R)$ --- we can write homomorphisms in the way we're used to in linear algebra, where $A \in \operatorname{Mat}_{m \times n}(R)$ sends $(x_1, \ldots, x_n)^t$ to $A(x_1, \ldots, x_n)^t$. So as long as we work with free modules and homomorphisms, to a large extent we can operate as if we're doing linear algebra. But in linear algebra, there's various characterizations of nondegenerate matrices that no longer hold here --- for instance, a linear operator that is injective (meaning it has zero kernel) is also surjective, but that's not true for general modules.

\subsection{Generators and Relations}

\vspace{0.75em}

\begin{definition}
    A collection of elements $m_1$, \ldots, $m_n \in M$ forms a system of \textbf{generators} if every $x \in M$ can be expressed as $\sum r_im_i$ for $r_i \in R$. 
\end{definition}

So in other words, $\varphi_{m_1, \ldots, m_n} : R^n \to M$ is onto. If such a finite set exists, we say that $M$ is \emph{finitely generated}. Many modules we're interested in are in fact finitely generated. 

If this map is also one-to-one, then it's an isomorphism, and $M$ is free. But usually this won't happen, and we still want to describe $M$. To do this, we can look at $K = \ker(\varphi)$, which is a submodule in $R^n$. If $K$ is itself finitely generated, then we can choose a set of generators for $K$, and get a somewhat explicit description of $M$ --- we can fix a system of $k$ generators for $K$, and obtain another homomorphism $\psi : R^k \to K$. Since $K$ sits inside $R^n$, we can think of $\psi$ instead as a homomorphism $\psi : R^k \to R^n$, with image $K$. But such a homomorphism corresponds to a matrix $A \in \operatorname{Mat}_{k\times n}$, and we have $M = R^n/AR^n$. 

\begin{definition}
If we can find a finite set of generators for $M$ such that the kernel is also finitely generated, then $M$ is called \textbf{finitely presented}. 
\end{definition}

We'll see that for many rings, the finitely presented modules are a very large class of modules --- in fact, for many rings, any finitely generated module is finitely presented. We'll also see how to make this description of a module very explicit when the ring is a Euclidean domain. In particular, taking the Euclidean domain to be $\ZZ$ will give us a classification of all finitely generated abelian groups, and taking the Euclidean domain to be $F[x]$ will give us Jordan normal form!
